% Public academic CV. Build from the repository root with:
% tectonic --outdir files _cv/Constantin_Auga_CV.tex
\documentclass[10pt,a4paper]{article}
\usepackage[left=17.5mm,right=17.5mm,top=13mm,bottom=13mm]{geometry}
\usepackage{libertine}
\usepackage{fontawesome5}
\usepackage{xcolor}
\usepackage{enumitem}
\usepackage{tabularx}
\usepackage{hyperref}
\definecolor{ink}{HTML}{202D3E}
\definecolor{accent}{HTML}{246781}
\definecolor{muted}{HTML}{52677E}
\definecolor{rulecolor}{HTML}{B5C8D4}
\hypersetup{
  colorlinks=true,urlcolor=accent,linkcolor=accent,
  pdftitle={Constantin Alexander Auga Academic Research CV},
  pdfauthor={Constantin Alexander Auga},
  pdfsubject={Research in computer vision and multimodal learning}
}
\pagestyle{empty}
\setlength{\parindent}{0pt}
\setlength{\parskip}{0pt}
\setlength{\tabcolsep}{0pt}
\hyphenpenalty=10000
\exhyphenpenalty=10000
\emergencystretch=1em
\setlist[itemize]{leftmargin=11pt,label=\textcolor{accent}{\scriptsize\textbullet},
  topsep=1.5pt,itemsep=0.8pt,parsep=0pt,partopsep=0pt}
\newcommand{\cvsection}[2]{%
  \par\vspace{6pt}\noindent
  {\color{accent}#1}\hspace{4pt}{\bfseries\fontsize{10.8}{12.2}\selectfont #2}\par
  \vspace{2pt}{\color{rulecolor}\hrule height 0.35pt}\vspace{4pt}
}
\newcommand{\cvrow}[2]{%
  \noindent\begin{tabularx}{\linewidth}{@{}Xr@{}}
    {\bfseries #1} & {\fontsize{9}{11}\selectfont\color{muted}#2}
  \end{tabularx}\par
}
\newcommand{\cvdetail}[1]{{\fontsize{9.3}{11}\selectfont\color{muted}\itshape #1}\par}
\newcommand{\cvlinks}[3]{\href{#1}{[Paper]}\enspace\href{#2}{[Project]}\enspace\href{#3}{[Code]}}
\begin{document}
\color{ink}
\fontsize{10.5}{12.2}\selectfont
\raggedright

\begin{center}
  {\bfseries\fontsize{23}{26}\selectfont Constantin Alexander Auga}\par
  \vspace{6pt}{\color{accent}\rule{30mm}{0.7pt}}\par
  \vspace{3pt}
  {\fontsize{9.3}{11.5}\selectfont
    {\color{accent}\faMapMarker*}\enspace Berlin, Germany\qquad
    {\color{accent}\faEnvelope}\enspace
    \href{mailto:constantin.auga@student.hpi.uni-potsdam.de}{constantin.auga@student.hpi.uni-potsdam.de}\par
    \vspace{3pt}
    \href{https://scholar.google.de/citations?hl=en\&user=INcDsjcAAAAJ}{\faGraduationCap\enspace Google Scholar}
    \quad{\color{rulecolor}|}\quad
    \href{https://github.com/Conscht}{\faGithub\enspace GitHub}
    \quad{\color{rulecolor}|}\quad
    \href{https://conscht.github.io/}{\faGlobe\enspace Website}
    \quad{\color{rulecolor}|}\quad
    \href{https://www.linkedin.com/in/constantin-alexander-auga-995a31230/}{\faLinkedin\enspace LinkedIn}
  }
\end{center}
\vspace{-7pt}

\cvsection{\faLightbulb}{Research Interests}
{\fontsize{10}{12.2}\selectfont
My research spans computer vision, audio, and multimodal learning, including generative modeling, self-supervised learning, and multimodal model evaluation.
I am particularly interested in how learned representations support perception, prediction, and world understanding, with a focus on multimodal video learning and dynamic environments.\par}

\cvsection{\faGraduationCap}{Education}
\cvrow{Hasso Plattner Institute, University of Potsdam}{Oct 2024 -- Jun 2027 (expected)}
M.Sc. IT Systems Engineering\quad{\color{rulecolor}|}\quad Current grade: \textbf{1.1} (1.0 best)\par
\vspace{3pt}
\cvrow{University of Technology Sydney}{Feb -- Jun 2027 (upcoming)}
Selected for an exchange semester\par
\vspace{3pt}
\cvrow{University of Hamburg}{Oct 2021 -- Sep 2024}
B.Sc. Computer Science\quad{\color{rulecolor}|}\quad Final grade: 2.0 (top 26\% of cohort)

\cvsection{\faBookOpen}{Publications and Preprints}
{\bfseries Beyond Accuracy: Benchmarking Cross-Task Consistency in Unified Multimodal Models}\par
\cvrow{WACV 2027 -- Accepted}{\cvlinks{https://arxiv.org/abs/2604.25072}{https://weixingw.github.io/xtc-bench/}{https://github.com/liudvikaszekas/xtc-benchmark}}
{\fontsize{9.1}{10.5}\selectfont
  \textbf{Constantin Alexander Auga}\textsuperscript{*}, Weixing Wang\textsuperscript{*},
  Liudvikas Zekas\textsuperscript{*}, Anton Hackl\textsuperscript{*}, et al.\par
  {\color{muted}\textsuperscript{*}Equal contribution. Co-first authors listed with the applicant first.}\par
}
\begin{itemize}
  \item Developed XTC-Bench, a scene-graph benchmark with 31,530 facts across 2,000 images.
  \item Tested whether unified representations yield consistent generation and understanding capabilities across nine models, including Gemini 2.5 Flash Image.
\end{itemize}
\vspace{3pt}
{\bfseries TargetSEC: Plug-and-Play In-the-Wild Speech Emotion Conversion via Arousal-Conditioned Latent Style Diffusion}\par
\cvrow{Under review at ICASSP 2027}{\cvlinks{https://arxiv.org/abs/2606.07293}{https://conscht.github.io/TargetSEC/}{https://github.com/Conscht/TargetSEC}}
{\fontsize{9.1}{10.5}\selectfont\textbf{Constantin Alexander Auga}}\par
\begin{itemize}
  \item Developed a latent diffusion model for speech emotion conversion with continuous arousal control.
  \item Outperformed non-duration baselines on MSP-Podcast in conversion accuracy and speech naturalness.
\end{itemize}

\cvsection{\faFlask}{Research Experience}
\cvrow{Hearing Beyond the Frame}{Master's thesis (ongoing)}
\cvdetail{Hasso Plattner Institute, University of Potsdam\quad |\quad Supervisor: Gerard de Melo}
\begin{itemize}
  \item Investigating when audio provides information beyond vision for predicting future actions.
  \item Developed a JEPA-inspired self-supervised feature predictor, achieving performance comparable to similarly sized supervised models.
\end{itemize}
\vspace{3pt}
\cvrow{Research Assistant -- Max Delbr\"uck Center}{Sep 2025 -- Present}
\cvdetail{Continuous convolutional neural networks (CCNNs)}
\begin{itemize}
  \item Evaluated resolution robustness of continuous convolutional networks on 2D and 3D data.
  \item Contributed to \href{https://www.aqqua.life/}{aQQua}, a tool for inspecting and interpreting learned embedding spaces.
\end{itemize}
\vspace{3pt}
\cvrow{Research Assistant -- German Aerospace Center (DLR)}{Sep 2023 -- Sep 2024}
\cvdetail{Machine learning for hydrogen-powered aviation}
\begin{itemize}
  \item Built real-time sensor pipelines for ML-based hydrogen-leak localization.
  \item Visualized spatial leak estimates as heatmaps for experimental analysis.
\end{itemize}

\cvsection{\faTrophy}{Awards and Academic Service}
\cvrow{1st Place, Jaguar Re-Identification Challenge (Kaggle)}{Mar 2026}
Ranked 1st of 370+ competitors with a DINOv3-based visual re-identification system.\par
\vspace{3pt}
\cvrow{Student Representative, Hasso Plattner Institute}{Nov 2025 -- Present}
Served on official selection committees for two professorships.

\cvsection{\faCode}{Technical Skills and Languages}
\textbf{Programming:} Python, C++, C\#, Java, JavaScript, SQL\par
\textbf{Machine learning:} PyTorch, Hugging Face, NumPy, pandas, scikit-learn\par
\textbf{Tools:} Git, Linux, ROS, Docker (basic familiarity)\par
\textbf{Languages:} German (native), English (TOEFL iBT 103, CEFR C1), Portuguese, Spanish, Georgian (basic)

\end{document}
